\begin{examplebox}
	\textbf{\textcolor{CExample}{例 1（基础）}}
	\textbf{\textcolor{CExample}{题目：}} 求函数 $y=2\sin(3x-\frac{\pi}{4})$ 的周期和值域。\textbf{\textcolor{CExample}{解题步骤：}}
	\begin{description}[style=nextline, leftmargin=2.8em, labelsep=0.5em, itemsep=0.45em]
		\item[\textbf{第一步（周期计算）：}]
			利用正弦型函数周期公式 $T=\frac{2\pi}{|\omega|}$，此处 $\omega=3$，故周期 $T=\frac{2\pi}{3}$。
			\par\textbf{理由：} 标准正弦函数周期为 $2\pi$，$\omega$ 压缩横坐标。
		\item[\textbf{第二步（值域计算）：}]
			振幅 $A=2$，所以值域为 $[-2,2]$。
			\par\textbf{理由：} $\sin$ 值域为 $[-1,1]$，乘以 $2$ 后得到。
	\end{description}
	\textbf{\textcolor{CExample}{关键结论：}}
	\textbf{周期 $T=\dfrac{2\pi}{3}$，值域 $[-2,2]$。}

	\medskip
	\textbf{\textcolor{CExample}{例 2（稍变形）}}
	\textbf{\textcolor{CExample}{题目：}} 判断函数 $f(x)=\cos(2x+\frac{\pi}{2})$ 的奇偶性，并求其在 $[0,\pi]$ 上的单调增区间。\textbf{\textcolor{CExample}{解题步骤：}}
	\begin{description}[style=nextline, leftmargin=2.8em, labelsep=0.5em, itemsep=0.45em]
		\item[\textbf{第一步（化简）：}]
			利用诱导公式 $\cos(\alpha+\frac{\pi}{2})=-\sin\alpha$，得 $f(x)=-\sin(2x)$。
			\par\textbf{理由：} $\cos(\alpha+\frac{\pi}{2})=-\sin\alpha$ 是常用诱导公式。
		\item[\textbf{第二步（奇偶性）：}]
			因为
			\[
				\begin{aligned}
					f(-x) &= -\sin(-2x) \\ &= \sin(2x) \\ &= -f(x),
			\end{aligned}
			\]
			所以 $f(x)$ 是奇函数。
			\par\textbf{理由：} 代入 $-x$ 并比较。
		\item[\textbf{第三步（单调增区间）：}]
			由 $y=-\sin(2x)$，考虑 $\sin(2x)$ 的减区间即是 $f(x)$ 的增区间。$\sin(2x)$ 在 $2x\in[-\frac{\pi}{2}+2k\pi,\frac{\pi}{2}+2k\pi]$ 上递增，则 $f(x)$ 在对应区间递减，故 $f(x)$ 的增区间满足 $2x\in[\frac{\pi}{2}+2k\pi,\frac{3\pi}{2}+2k\pi]$，即 $x\in[\frac{\pi}{4}+k\pi,\frac{3\pi}{4}+k\pi]$。取 $k=0$ 得在 $[0,\pi]$ 上的增区间为 $[\frac{\pi}{4},\frac{3\pi}{4}]$。
			\par\textbf{理由：} 负号改变单调性，复合函数求单调性。
	\end{description}
	\textbf{\textcolor{CExample}{关键结论：}}
	\textbf{奇函数，在 $[0,\pi]$ 上增区间为 $[\frac{\pi}{4},\frac{3\pi}{4}]$。}

	\medskip
	\textbf{\textcolor{CExample}{例 3（综合应用）}}
	\textbf{\textcolor{CExample}{题目：}} 求函数 $y=\tan(2x+\frac{\pi}{3})$ 的定义域和周期。\textbf{\textcolor{CExample}{解题步骤：}}
	\begin{description}[style=nextline, leftmargin=2.8em, labelsep=0.5em, itemsep=0.45em]
		\item[\textbf{第一步（周期计算）：}]
			正切型函数周期公式 $T=\frac{\pi}{|\omega|}$，$\omega=2$，所以 $T=\frac{\pi}{2}$。
			\par\textbf{理由：} 正切标准周期为 $\pi$，$\omega$ 压缩。
		\item[\textbf{第二步（定义域求解）：}]
			要求 $2x+\frac{\pi}{3}\neq\frac{\pi}{2}+k\pi$，解得 $x\neq\frac{\pi}{12}+\frac{k\pi}{2}$，$k\in\mathbb{Z}$。
			\par\textbf{理由：} 正切函数在 $\frac{\pi}{2}+k\pi$ 处无定义。
	\end{description}
	\textbf{\textcolor{CExample}{关键结论：}}
	\textbf{周期 $T=\dfrac{\pi}{2}$，定义域为 $\{x\mid x\neq\dfrac{\pi}{12}+\dfrac{k\pi}{2},k\in\mathbb{Z}\}$。}
\end{examplebox}
